\documentclass[12pt, a4paper]{article}
\usepackage{luatexja}
\usepackage{luatexja-fontspec}
\usepackage{fontspec}
\usepackage{geometry}
\usepackage{booktabs}
\usepackage{array}
\usepackage{hyperref}
\usepackage{enumitem}
\usepackage{titlesec}
\usepackage{parskip}
\usepackage{longtable}
\usepackage{float}
\usepackage{amsmath}
\usepackage{xcolor}
\usepackage{tcolorbox}
\usepackage{mdframed}
\usepackage{multirow}
\usepackage{tabularx}

\geometry{top=2.5cm, bottom=2.5cm, left=3cm, right=3cm}
\setmainfont{Times New Roman}
\setmainjfont[BoldFont={Microsoft YaHei}]{MingLiU_HKSCS}

\titleformat{\section}{\large\bfseries}{}{0em}{}[\titlerule]
\titleformat{\subsection}{\normalsize\bfseries}{}{0em}{}

\hypersetup{colorlinks=true, linkcolor=blue, urlcolor=blue}

\newcommand{\done}{\textcolor{green!50!black}{$\checkmark$}}
\newcommand{\inprog}{\textcolor{orange!70!black}{進行中}}
\newcommand{\todo}{\textcolor{gray}{待完成}}

\begin{document}

\begin{center}
  {\LARGE\bfseries MangaSim：漫畫盜版隱性偏好的多Agent AI模擬}\\[0.4em]
  {\large 作業四進度報告}\\[0.6em]
  涂皓 (tuhow)\quad 花曉桐\quad 徐瑋倢\quad 黃少桐\\[0.3em]
  2026年5月
\end{center}

\vspace{0.5em}
\noindent\rule{\linewidth}{0.5pt}

% ─── 1. Background ────────────────────────────────────────────────────────────
\section{1. 研究背景與問題}

\subsection{1.1 為什麼選漫畫？}

本研究以\textbf{漫畫盜版}作為研究核心，理由如下：

\begin{itemize}
  \item \textbf{盜版規模最大}：依據MUSO（2024）統計，漫畫佔全球出版類盜版瀏覽量的
        70.15\%，2024年全球漫畫盜版瀏覽量達\textbf{466億次}，是唯一持續成長的盜版類別。\footnote{MUSO, ``2024 Manga Piracy Trends and Insights,'' \url{https://www.muso.com/muso-2024-manga-piracy-trends-and-insights-report}}

  \item \textbf{台灣脈絡強}：2024年文化消費調查顯示，台灣漫畫讀者中\textbf{66.7\%使用數位閱讀}；
        日本漫畫在台灣授權常有空窗期與繁中版時差，形成明顯的「合法無法取得」誘因。

  \item \textbf{社會期許偏誤（SDB）效果明顯}：漫畫盜版行為有特定平台可指稱
        （如Bato.to、MangaDex），受訪者在具名情境下較容易自我審查，有助量化SDB差距。
\end{itemize}

\subsection{1.2 研究問題}

\begin{enumerate}
  \item 不同社會經濟背景的台灣漫畫讀者，其盜版使用頻率有何差異？
  \item 匿名與具名情境下，漫畫盜版行為的自我揭露有多大落差？
  \item 價格、可得性、便利性、時差等因素，對各群體的影響力如何排序？
  \item 若推出完整正版繁中漫畫平台，各群體的付費意願（WTP）為何？
\end{enumerate}

\subsection{1.3 台灣讀者常用盜版平台（2023–2025）}

下表為Gemini深度調查報告（綜合Similarweb、Semrush流量數據）整理的八大台灣核心盜版平台：

\begin{table}[H]
\centering\small
\begin{tabular}{p{2.3cm}p{2.2cm}p{1.6cm}p{1.8cm}p{4.3cm}}
\toprule
\textbf{站名} & \textbf{主網址} & \textbf{月流量} & \textbf{台灣佔比} & \textbf{特色與現況} \\
\midrule
Manhuagui\\（漫畫櫃） & manhuagui.com\\tw.manhuagui.com & 6,489萬 & 77.06\% & 台灣第一大盜版漫畫站；設有繁中專屬子網域；平均停留21分鐘、每次13.8頁；穩定運作 \\
\midrule
Happymh\\（嗨皮漫畫） & happymh.com & 5,244萬 & 67.31\% & 行動端最強；台灣社群網站分類第8名；原生支援繁簡自動切換；平均停留15分鐘 \\
\midrule
MangaDex\\（繁中版） & mangadex.org & 3.12億（全球） & 繁中主力 & 社群驅動、無廣告、API開放；台灣本土漢化組聚集地；提供高品質繁中社群譯 \\
\midrule
Baozimh\\（包子漫畫） & baozimh.com & 278萬 & 53.58\% & 專攻韓漫條漫；行動端87\%；網域頻繁更換 \\
\midrule
Manhuaren\\（漫畫人） & manhuaren.com & 283萬 & 60.29\% & 繁中支援含港台專區；高度依賴非官方APK散布 \\
\midrule
DM5\\（動漫屋） & dm5.com & 約200萬 & 39.34\% & 老牌站；繁簡切換；商業廣告含加密貨幣聯盟 \\
\midrule
Copymanga\\（拷貝漫畫） & mangacopy.com & 歷史極高 & 不明 & 主網域多次遭封；靠社群（PTT、Dcard）傳播備用連結；高度不穩定 \\
\midrule
18comic\\（禁漫天堂） & 18comic.org & 28.3萬 & 54.28\% & 原生繁中；成人內容為主；台灣ISP頻繁封鎖 \\
\bottomrule
\end{tabular}
\caption{台灣讀者八大核心漫畫盜版平台（2023–2025）。來源：Gemini深度調查報告，綜合Similarweb、Semrush}
\end{table}

\textbf{重要背景：2024年iWIN事件}——台灣iWIN機構以「二次元虛擬兒少」為由，向PTT、Dcard等合法平台發送大規模下架通知，引發ACG社群強烈反彈。此事件意外推動大批台灣讀者轉向境外盜版平台，使Manhuagui等站流量進一步集中。

\medskip
\noindent\textbf{台灣讀者常用合法平台}：ReadMoo（讀墨）、BookWalker台灣、Comico、LINE Webtoon、MANGA Plus（集英社免費）

% ─── 2. Survey Method ─────────────────────────────────────────────────────────
\section{2. 問卷設計方法}

\subsection{2.1 回應蒐集方法：SSR vs DLR}

本研究採用\textbf{語義相似度評分（Semantic Similarity Rating, SSR）}蒐集模擬回應，
以取代傳統直接Likert評分法（Direct Likert Response, DLR）。

\textbf{DLR的問題}：直接要求LLM輸出1–5的整數，會導致嚴重的中央集中偏誤，
回應集中在2–4，無法區分細微態度差異。
依據Maier等人（2025）的驗證，DLR與真實人類分布的KS相似度僅0.26。\footnote{Benjamin F. Maier et al., ``LLMs Reproduce Human Purchase Intent via Semantic Similarity Elicitation of Likert Ratings,'' arXiv:2510.08338 (2025). 以9,300份人類問卷、57項消費品調查驗證；SSR達人類重測信度的90\%，KS相似度大於0.85。}

\medskip
\textbf{SSR流程}：
\begin{enumerate}[leftmargin=2em]
  \item 向LLM提問，\textbf{要求自由文字回應}（不指定輸出格式）
  \item 對LLM回應進行\textbf{文字嵌入}（OpenAI \texttt{text-embedding-3-small}）
  \item 同樣對\textbf{5個描述性錨點陳述}進行嵌入
  \item 計算回應與各錨點的\textbf{餘弦相似度}，轉換為概率分布（pmf）
  \item 取期望值作為Likert分數（可得非整數，分布更自然）
\end{enumerate}

SSR的優點：LLM以自然語言回答，不受數字格式限制；錨點設計確保量表一致性。

\subsection{2.2 社會期許偏誤（SDB）的處理}

\textbf{SDB的問題}：受訪者在具名情境下，對帶有道德污名的行為（如使用盜版）會傾向低報。

\textbf{本研究的解法}：在問卷中加入\textbf{匿名/具名情境配對題}，
對同一個模擬受訪者（persona）分別在兩種情境下提問，
直接計算\textbf{SDB差距}（匿名均分 $-$ 具名均分）。
差距為正值代表匿名時更願意承認使用盜版。

\begin{center}
\textbf{SDB差距} $=$ 匿名情境回應均值 $-$ 具名情境回應均值
\end{center}

% ─── 3. Survey Questions ──────────────────────────────────────────────────────
\section{3. 問卷題目設計}

\subsection{3.1 題目總覽}

問卷共14題，包含：4題配對SDB核心題（2對匿名/具名）、1題多選驅動因素、
7題態度量表、1題付費意願、1題開放式。

\begin{longtable}{p{2cm}p{3cm}p{8.5cm}p{1.5cm}}
\toprule
\textbf{題號} & \textbf{主題} & \textbf{題目文字} & \textbf{方法} \\
\midrule
\endhead
Q1a-匿 & 線上閱讀頻率（匿名）
  & 【匿名情境】過去3個月，請問你在漫畫櫃（Manhuagui）、嗨皮漫畫（Happymh）、MangaDex、拷貝漫畫等非官方網站閱讀漫畫的頻率？
  & SSR / 匿名 \\[4pt]
Q1a-名 & 線上閱讀頻率（具名）
  & 【具名情境】過去3個月，請問你在漫畫櫃（Manhuagui）、嗨皮漫畫（Happymh）、MangaDex、拷貝漫畫等非官方網站閱讀漫畫的頻率？
  & SSR / 具名 \\[4pt]
Q1b-匿 & 下載頻率（匿名）
  & 【匿名情境】過去3個月，請問你從非官方管道下載漫畫圖檔（如PDF、CBZ）儲存在手機或電腦的頻率？
  & SSR / 匿名 \\[4pt]
Q1b-名 & 下載頻率（具名）
  & 【具名情境】過去3個月，請問你從非官方管道下載漫畫圖檔儲存在手機或電腦的頻率？
  & SSR / 具名 \\[4pt]
Q2 & 合法平台使用率
  & 過去3個月，請問你使用ReadMoo（讀墨）、BookWalker、Comico、LINE Webtoon、MANGA Plus等官方平台閱讀漫畫的頻率？
  & SSR \\[4pt]
Q3 & 使用非官方站的原因
  & 你選擇漫畫櫃、嗨皮漫畫等非官方網站的原因有哪些？（可複選）
  & 多選 \\[4pt]
Q4 & 道德評價
  & 整體來說，你認為使用漫畫櫃、嗨皮漫畫等非官方網站閱讀漫畫有多不道德？
  & SSR \\[4pt]
Q5 & 版權制度公平性
  & 你認為台灣目前的漫畫授權與版權制度，對一般讀者來說是否公平？
  & SSR \\[4pt]
Q6 & 「沒辦法才用」的接受度
  & 當你想看的漫畫沒有合法繁中版，或者合法版本要貴很多時，你覺得使用非官方網站可以接受嗎？
  & SSR \\[4pt]
Q7 & 罪惡感
  & 使用漫畫櫃或嗨皮漫畫等非官方網站看漫畫時，你心裡有多少罪惡感？
  & SSR \\[4pt]
Q8 & 合法平台CP值
  & 以ReadMoo或BookWalker台灣版的定價，你覺得是否值得？
  & SSR \\[4pt]
Q9 & 監管態度
  & 你認為政府對非官方漫畫網站的管制（例如DNS封鎖、要求下架等），執行力度是否適當？
  & SSR \\[4pt]
WTP & 付費意願
  & 假設有一個收錄完整日漫繁中版的官方訂閱平台，你最多願意每月付多少台幣？不願付請填0。
  & WTP \\[4pt]
Q-開 & 開放回應
  & 請用一到兩句話說說你對台灣漫畫非官方網站問題的看法，或你繼續使用或不使用的原因。
  & 開放式 \\
\bottomrule
\caption{MangaSim問卷題目清單（共14題）}
\end{longtable}

\noindent
\textbf{SDB差距A（線上閱讀）} $=$ Q1a-匿均值 $-$ Q1a-名均值\\
\textbf{SDB差距B（下載行為）} $=$ Q1b-匿均值 $-$ Q1b-名均值

\subsection{3.2 頻率題SSR錨點設計}

Q1a與Q1b共用以下5個行為描述錨點，\textbf{不使用數字或模糊頻率詞}：

\begin{table}[H]
\centering\small
\begin{tabular}{cp{10.5cm}p{2cm}}
\toprule
\textbf{分數} & \textbf{錨點陳述（行為描述）} & \textbf{大約頻率} \\
\midrule
1 & 「過去3個月我完全沒有用過這類網站或下載過掃描檔。」 & 0次 \\
2 & 「偶爾用過，大概只有1到3次，不算是習慣。」 & 不到每月一次 \\
3 & 「某些系列我每週會去看一兩次更新。」 & 每週1至2次 \\
4 & 「幾乎每天都會開這類網站追更新。」 & 接近每天 \\
5 & 「這是我過去3個月看漫畫的主要方式，大部分內容都從這裡讀。」 & 主要管道 \\
\bottomrule
\end{tabular}
\caption{頻率題SSR錨點（Q1a / Q1b共用）}
\end{table}

\subsection{3.3 Q3多選題選項設計}

\begin{table}[H]
\centering\small
\begin{tabular}{lp{11cm}}
\toprule
\textbf{選項代碼} & \textbf{選項內容} \\
\midrule
D1 & 官方平台定價太貴（ReadMoo單冊售價、BookWalker訂閱費偏高） \\
D2 & 想看的漫畫在台灣沒有合法繁體中文版（日本原版未授權給台灣） \\
D3 & 官方繁中版更新太慢——非官方翻譯往往在日本原版發布後數小時內就上線，比官方早數週甚至數月 \\
D4 & 使用方便：漫畫櫃、嗨皮漫畫等網站畫質好、不強制廣告、可手機閱讀 \\
D5 & 習慣問題：從國中或高中就開始用了，沒有特別想換 \\
D6 & 覺得對實際創作者的傷害有限（尤其是大型日本出版社） \\
D7 & 因為iWIN事件等審查爭議，不信任台灣的內容管理方式，主動選擇境外平台 \\
\bottomrule
\end{tabular}
\caption{Q3驅動因素選項（D3「時差」與D7「iWIN監管不信任」為台灣特有驅動因素）}
\end{table}

% ─── 4. Persona Design ────────────────────────────────────────────────────────
\section{4. 模擬受訪者設計（Persona）}

\subsection{4.1 設計邏輯}

四個模擬受訪者依據三個層次設計：
\begin{itemize}
  \item \textbf{人口結構}：對照2024文化調查漫畫讀者段（N=693），涵蓋台灣主要閱讀人口分布。
  \item \textbf{研究文獻}：Phau \& Liang（2012）\footnote{Phau \& Liang, ``Downloading Digital Video Games: Predictors, Moderators, and Consequences,'' \textit{Marketing Intelligence and Planning} 30(7) (2012). 原研究對象為電玩；本研究類推至漫畫，同為日本IP密集的數位娛樂產品。} 指出低收入學生族群的數位盜版頻率最高，SDB差距也最顯著。
  \item \textbf{SDB理論}：四個受訪者在「道德顧慮」與「習慣性使用非官方網站」兩軸上刻意差異化，預測SDB差距排序：學生 $>$ 一般上班族 $>$ 技術工作者 $>$ 注重正版的家長。
\end{itemize}

\subsection{4.2 受訪者設定}

\begin{table}[H]
\centering\small
\begin{tabular}{p{2.8cm}cccp{5.5cm}}
\toprule
\textbf{受訪者} & \textbf{年齡} & \textbf{月收入} & \textbf{居住地} & \textbf{漫畫閱讀習慣} \\
\midrule
Manga Otaku Student
  & 20歲 & NT\$8,000 & 台北
  & 每天手機用漫畫櫃追20部以上；對大出版商無罪惡感；具名壓力最大 \\
Casual Pirate Worker
  & 28歲 & NT\$35,000 & 高雄
  & 通勤時用嗨皮漫畫看5至6部；高中養成習慣；偶買實體版；正式場合不承認 \\
Legal-Leaning Parent
  & 42歲 & NT\$80,000 & 台中
  & 偶爾看；付費ReadMoo和Comico；關注iWIN事件；反感盜版站廣告環境 \\
Tech Pro Subscriber
  & 32歲 & NT\$120,000 & 台北
  & 訂閱BookWalker日本版；偶用MangaDex高品質繁中社群譯；不用Manhuagui機轉版 \\
\bottomrule
\end{tabular}
\caption{模擬受訪者設計（漫畫讀者）}
\end{table}

\subsection{4.3 完整模擬提示詞}

\begin{mdframed}[backgroundcolor=gray!8, linecolor=gray!40, linewidth=0.5pt]
\textbf{Persona 1 --- Manga Otaku Student}\\[0.3em]
\small
You are a 20-year-old university student in Taipei living on NT\$8,000/month.
You follow over 20 ongoing manga series and read every day, almost entirely on your phone.
Your main platforms are Manhuagui (漫畫櫃, tw.manhuagui.com) for Traditional Chinese scanlations,
and MangaDex for series where local fan translation groups post higher-quality versions.
You cannot afford subscriptions --- ReadMoo has a limited catalogue and charges per volume,
which adds up fast. Official Traditional Chinese releases often lag weeks or months behind
the Japanese originals, and you have no patience for that. You feel almost no moral guilt:
you see the big Japanese publishers as distant corporations, not creators who suffer from
your reading habits. However, you would not openly admit to using these sites in front of
professors, employers, or on a named survey --- it just sounds bad.
\end{mdframed}

\begin{mdframed}[backgroundcolor=gray!8, linecolor=gray!40, linewidth=0.5pt]
\textbf{Persona 2 --- Casual Pirate Worker}\\[0.3em]
\small
You are a 28-year-old office worker in Kaohsiung earning NT\$35,000/month.
You casually follow five or six manga series. You read them on Happymh (嗨皮漫畫)
on your phone during commutes --- a habit you picked up in high school that you never
bothered to change. You buy physical collector's editions of your favorite series once
in a while, which makes you feel like you are supporting the creators enough.
You know using these sites is technically copyright infringement but it does not feel
wrong to you. In a formal or professionally named context, you would deny or downplay
your use of piracy sites, because it could reflect badly on you.
\end{mdframed}

\begin{mdframed}[backgroundcolor=gray!8, linecolor=gray!40, linewidth=0.5pt]
\textbf{Persona 3 --- Legal-Leaning Parent}\\[0.3em]
\small
You are a 42-year-old parent of two in Taichung with a household income of NT\$80,000/month.
You read manga occasionally --- mostly titles your teenage children recommend.
You pay for ReadMoo and Comico because you want to set a good example for your children about
respecting intellectual property. You find sites like Manhuagui cluttered, full of sketchy ads,
and you do not trust them. The 2024 iWIN content moderation controversy made you more aware
of the risks of unregulated platforms. You feel clearly uncomfortable about copyright
infringement and would not use piracy sites even if a title is not available legally.
\end{mdframed}

\begin{mdframed}[backgroundcolor=gray!8, linecolor=gray!40, linewidth=0.5pt]
\textbf{Persona 4 --- Tech Pro Subscriber}\\[0.3em]
\small
You are a 32-year-old software engineer in Taipei earning NT\$120,000/month.
You subscribe to BookWalker Japan for simulpublication of series you care about.
You also use MangaDex occasionally for fan translations of niche titles not yet officially
licensed in Taiwan --- you appreciate MangaDex's no-ad, API-based design and the quality
of its community translators over the machine-converted Traditional Chinese on Manhuagui.
You rationalize MangaDex use as filling a gap the market created, not as theft.
You would pay comfortably for a well-designed all-in-one Traditional Chinese manga platform
if it existed. In professional settings, you would not mention using any unofficial sites.
\end{mdframed}

\subsection{4.4 與2024文化調查的對應}

\begin{table}[H]
\centering\small
\begin{tabular}{p{4.5cm}rp{5.5cm}}
\toprule
\textbf{調查指標（漫畫讀者段，N=693）} & \textbf{數值} & \textbf{對應受訪者} \\
\midrule
數位閱讀使用率 & 66.7\% & 全部（主要使用數位） \\
實體閱讀使用率 & 58.4\% & Casual Worker（買實體珍藏版） \\
閱讀頁漫（翻頁式）& 72.4\% & Otaku Student、Casual Worker \\
閱讀條漫（長條式）& 37.8\% & Tech Pro（LINE Webtoon） \\
圖書館紙本借閱漫畫 & 20.1\% & Legal Parent（合法途徑） \\
閱讀台灣創作者漫畫 & 19.0\% & 主要集中在Legal Parent \\
\bottomrule
\end{tabular}
\caption{2024文化調查漫畫讀者段與模擬受訪者設定的對應關係}
\end{table}

目前四個受訪者各佔25\%為暫定值；正式執行前將依調查中年齡$\times$收入分布調整權重。

% ─── 5. Pilot Results ─────────────────────────────────────────────────────────
\section{5. 初步測試結果（DLR方法基線）}

\begin{tcolorbox}[colback=yellow!10, colframe=orange!50, boxrule=0.5pt]
以下結果來自先前以\textbf{DLR方法}（直接Likert評分）進行的廣泛盜版問卷初步測試。
受中央集中偏誤影響，\textbf{絕對數值僅供參考；排名方向性可信，將以SSR重新驗證。}
\end{tcolorbox}

\subsection{5.1 SDB差距排行}

\begin{table}[H]
\centering\small
\begin{tabular}{lrp{8cm}}
\toprule
\textbf{受訪者} & \textbf{SDB差距} & \textbf{說明} \\
\midrule
預算有限的學生族 & \textcolor{red!70!black}{\textbf{+3.20}} & 符合預期：使用頻率高、罪惡感低，但具名時社會壓力最強 \\
一般上班族 & \textcolor{red!70!black}{\textbf{+3.05}} & 習慣性使用加上職場形象顧慮，具名時顯著低報 \\
注重家庭的家長 & +1.65 & 本身使用頻率低，但具名時更強調道德立場 \\
高收入技術工作者 & +1.50 & 偶爾使用、收入較高，具名情境調整幅度最小 \\
\bottomrule
\end{tabular}
\caption{初步SDB差距排行（DLR基線，N=20）。低收入高使用族群的差距最大，符合理論預測。}
\end{table}

排名順序\textbf{符合理論預期}，驗證受訪者設計邏輯的有效性。
正式漫畫問卷以SSR執行後，預期Manga Otaku Student與Casual Worker排名最前。

\subsection{5.2 付費意願基線}

\begin{table}[H]
\centering\small
\begin{tabular}{lrp{6cm}}
\toprule
\textbf{受訪者} & \textbf{平均WTP（NT\$/月）} & \textbf{備註} \\
\midrule
預算有限的學生族 & NT\$352 & 約一個Netflix基本方案 \\
一般上班族 & NT\$268 & 低於學生（更強烈的反彈心理） \\
注重家庭的家長 & NT\$793 & 願意為家庭多付費 \\
高收入技術工作者 & NT\$1,000+ & 符合高收入特徵 \\
\bottomrule
\end{tabular}
\caption{付費意願基線（DLR，廣泛盜版問卷）。漫畫專題WTP預期更低（範疇縮小）。}
\end{table}

% ─── 6. Next Steps ────────────────────────────────────────────────────────────
\section{6. 下一步計畫}

\begin{enumerate}
  \item \textbf{5/15前（第一階段檢核點）}
    \begin{itemize}
      \item 確認漫畫問卷最終版（已含SSR錨點，見第3節）
      \item 完成SSR流程實作：向LLM收自由文字 $\rightarrow$ embedding $\rightarrow$ cosine similarity轉pmf
      \item 小規模漫畫問卷SSR測試（每組受訪者N=5，驗證流程）
    \end{itemize}
  \item \textbf{5/15至5/31}
    \begin{itemize}
      \item 正式執行漫畫問卷SSR，每組N$\geq$30
      \item 對比實驗：相同問卷跑DLR與SSR，比較SDB差距的分布差異
      \item 依2024調查漫畫讀者段資料調整受訪者權重
    \end{itemize}
  \item \textbf{期末前}
    \begin{itemize}
      \item 統計檢定：ANOVA、單調性檢定、排名相關係數
      \item 與2024調查比對：AI模擬的消費分布 vs 真實漫畫讀者分布
      \item 期末報告與簡報
    \end{itemize}
\end{enumerate}

\vspace{1em}
\noindent\rule{\linewidth}{0.5pt}
{\small
\textbf{資料來源：}
MUSO（2024年漫畫盜版報告）；
Similarweb流量數據（引自Variety, 2024）；
TorrentFreak盜版平台追蹤報導；
2024年文化內容消費趨勢調查（N=2,000，助教提供）；
Maier et al., arXiv:2510.08338 (2025)；
Phau \& Liang, \textit{Marketing Intelligence and Planning} 30(7) (2012)。
}

\end{document}
